\section{Preliminaries}


\subsection{DNS in a Nutshell} \label{subsec:dns}
The Domain Name System (DNS) is a decentralized database that translates human-readable domain names into  IP addresses. For example, when a user visits the website \name{www.sigplan.org}, the web browser  sends a DNS request to retrieve the IP address of the server hosting the site.

\paragraph{Namespace} 
The DNS namespace is organized hierarchically in a tree structure, 
with a root domain at the top, followed by top-level domains (TLDs), second-level domains (SLDs), and further subdomains. 
\Cref{fig:namespace} shows a small portion of this namespace, including the 
TLDs \name{org} and \name{com} as well as  SLDs such as 
\name{sigplan}. 
A domain name such as \name{www.sigplan.org} uniquely identifies a path in this tree, with the associated labels concatenated  bottom up, using the `.' character; the root node has an empty label, which is typically omitted. 
The entire namespace is divided into \emph{zones} or subtrees,
each managed by authoritative \emph{nameservers}.

\paragraph{Resource Records}
A DNS \emph{resource record} (RR) stores 
specific information about a domain name,
including the owner name (where this RR resides), record type, time to live (TTL), and 
value. In general, multiple
RRs with the same owner name and type can exist, provided they have
different values. 
Common types of RRs include 
\rtype{A}, \rtype{NS}, \rtype{SOA}, \rtype{TXT},  \rtype{CNAME}, and \rtype{DNAME}, 
each serving a distinct purpose. 
For example, an \rtype{A} record 
maps a domain name to an IPv4 address.
An \rtype{NS} record
specifies the authoritative nameservers for a domain.
A \rtype{CNAME} record
maps an alias domain name, e.g., \name{example.net}, 
to its canonical name, e.g., \name{www.example.com}; any request for \name{example.net} will be redirected to \name{www.example.com}.
A \rtype{DNAME} record
provides redirection for an \emph{entire subtree} of the domain name space to another domain.
For instance, 
a \rtype{DNAME} record at \name{legacy.example.com} with value \name{new.example.com}
would cause a query for \name{any.legacy.example.com} 
to be restarted for \name{any.new.example.com}.



A \emph{wildcard} record is
a special type (\name{$\star$}) of DNS record that
 matches requests for non-existent domain names. 
 This simplifies DNS management by allowing a single record to cover multiple potential queries.
For instance, suppose there is an \rtype{A} record
with \name{1.2.3.4} for \name{$\star$.example.org}.
Queries for  \name{one.example.org} or \name{any.example.org}
will all resolve to \name{1.2.3.4}, 
provided that no other records exist for these specific names.

\paragraph{Name Resolution}
A client  such as a web browser
initiates a DNS request by
directing the query to a \emph{recursive resolver} 
that handles the resolution process on the client's behalf.
If the recursive resolver has a cached response from a prior query, it can promptly reply to the client.
Otherwise,
it queries a hierarchy of authoritative nameservers to retrieve the answer.

\begin{figure}
\begin{minipage}[b]{.33\linewidth}
\centering
   \includegraphics[width=\textwidth]{img/namespace.pdf}
\captionsetup{skip=10pt}
   \caption{A small portion of the DNS namespace. 
   Dashed lines mark the zone cuts. }
   \label{fig:namespace}
\end{minipage}
\hfill
\begin{minipage}[b]{.65\linewidth}
\centering
   \includegraphics[width=\textwidth]{img/dns.pdf}
\captionsetup{skip=0pt}
   \caption{Name resolution for the type \rtype{A} query \name{www.sigplan.org}, with an empty cache at the DNS resolver. }
   \label{fig:dns}
\end{minipage}
\end{figure}

\begin{example}[Name Resolution] \label{example:resolution}
\Cref{fig:dns} depicts a DNS name resolution process. 
A client initiates a type \rtype{A} query for the name \name{www.sigplan.org} (\ding{182}). 
Assuming an empty cache, the recursive resolver forwards this query to one of the root nameservers  (\ding{183}). 
Here, the root nameserver cannot answer the query as it is only authoritative for the \name{root} zone. Hence, it refers the resolver to the authoritative nameservers for the \name{org} zone by  attaching the associated \rtype{NS} records and IP addresses in the referral (\ding{184}). 
The resolver then asks one of the \name{org} nameservers (\ding{185}). As the nameserver is not authoritative for the full domain,  it further refers the resolver to the authoritative nameservers for \name{sigplan.org} (\ding{186}). The resolver  finally receives the result, the IP address for \name{www.sigplan.org} (\ding{187} and \ding{188}), which it then sends back to the client (\ding{189}).
\end{example}

Note that, owing to  \emph{caching},
if the same query needs to be resolved again, the resolver can answer it directly without contacting the nameservers, provided the records' TTL has not expired.
Likewise, for a query of another name within an already cached zone, the resolver only needs to reach out to the relevant nameservers rather than starting from the root.
Note also that when a resolver has multiple nameservers to choose from, the selection is typically based on factors like response time and availability. DNS RFCs do not prescribe a specific nameserver selection algorithm, leaving this to the developers’ discretion.

\new{\paragraph{Subqueries}
During DNS resolution, the recursive resolver may need to perform additional queries to collect all the necessary information to resolve a client query. 
For instance,
suppose that a user wants to visit \name{example.org} and
 the authoritative
nameserver for the \name{example.org} zone is \name{sub.query.com},
the resolver must first resolve \name{sub.query.com}.\footnote{\new{This is a randomly chosen name just for illustrative purposes. As long as the authoritative
nameserver's name  differs from \name{example.org}, the resolver must then make additional subqueries to resolve it.}}
This process begins by querying the \name{com} TLD, followed by \name{query.com} (assuming the authoritative nameserver for \name{query.com} provides the IP address for \name{sub.query.com}). 
 These two additional queries are called \emph{subqueries}, as they are initiated by the resolver to resolve the original query.
}















\subsection{Labelled Transition Systems}\label{subsec:lts}
Given an alphabet $\Sigma$, a \emph{string} $str \in A^*$ is a sequence of labels $l_0 \cdot \ldots \cdot l_k$, where $l_0$ is the first label and $l_k$ is the last label. The length of the sequence, $k+1$, is denoted by $|str|$. For $i < j < |str|$, $str[i\ldots j]$ denotes the sequence $l_i\ldots l_j$ and $\epsilon$ denotes the \emph{empty} sequence. \new{Given two strings $str$ and $str'$, we write $str \preceq str'$ ($str \prec str'$) to denote that $str$ is a (proper) prefix of $str'$.} For notational convenience, we sometimes drop the explicit binary operator for concatenation.

\begin{definition}
	A \emph{labelled transition system} (LTS) \cite{heussner2012reachability}  $\cal A = \tup{S,S_{\mathit{init}},A,\trans}$ is given by a set of \emph{states} $S$, a set of initial states $S_{\mathit{init}}\subseteq S$, an alphabet $A$ of actions, and a set of transition rules $\scrsf{trans}$ of the form $s \xrightarrow{a} s'$, where $s,s'\in S$ and $a \in A$. As usual, we define an \emph{automaton} $\cal F = \tup{S, S_{\mathit{init}}, A, \trans, S_{\mathit{fin}}}$ as the LTS $\cal A$ along with a set of \emph{final} states $S_{\mathit{fin}} \subseteq S$. The automaton $\cal F$ is called a \emph{finite automaton} (FA) if the sets $S$ and $A$ are finite.
\end{definition}

A \emph{run} in the LTS $\cal A$ is a finite sequence $\rho = s_0 \xrightarrow{a_1} s_1 \ldots s_{n-1} \xrightarrow{a_n} s_n$, for $s_i \in S$ such that for every $1 \leq i \leq n$, $s_{i-1} \xrightarrow{a_i} s_i \in \trans$. We say that $\rho$ is a run from $s_0$ to $s_n$. The length of $\rho$ is $n$ and is denoted by $|\rho|$. The \emph{trace} of the run $\rho$ is the sequence of actions $a_1\ldots a_n$. Two runs are \emph{trace equivalent} if they have the same traces. Given the FA $\cal F$, a run in $\cal A$ is \emph{accepting} if it starts in an initial state and ends in a final state.

A state $s \in S$ is \emph{reachable} in $\cal A$ if there exists a run of $\cal A$ from an initial state $s_0$ to $s$.  The \emph{LTS verification} \emph{problem} asks whether, given an LTS $\cal A = \tup{S,S_{\mathit{init}},A,\trans}$, and an automaton $\cal F'$
over strings in $A^*$, there exists a run $\rho$ in $\cal A$ from an initial state, and an accepting, trace equivalent run $\rho'$ in $\cal F'$. We call the corresponding trace a \emph{witness} trace. The problem can be analogously defined for any system with an LTS semantics.






	









		












